\newpage
\section{Marco Teórico}

\vspace{0.5cm}

\noindent\textbf{1. Fundamentos de Ciberseguridad en la Educación Superior}
\begin{itemize}
    \item \textbf{1.1} Principios de la seguridad de la información (Confidencialidad, Integridad y Disponibilidad) (Stallings, 2020).
    \item \textbf{1.2} Amenazas y vulnerabilidades comunes en sistemas (Ingeniería social, phishing, malware) (Mitnick \& Simon, 2023).
    \item \textbf{1.3} Gestión y respuesta a incidentes cibernéticos (Fases de preparación, detección y contención) (Cichonski et al., 2022).
    \item \textbf{1.4} Desafíos actuales en la enseñanza puramente teórica de la seguridad informática.
\end{itemize}

\newpage

\noindent\textbf{2. Videojuegos Educativos (Serious Games)}
\begin{itemize}
    \item \textbf{2.1} Definición y características principales de los \textit{serious games} (Dörner et al., 2016).
    \item \textbf{2.2} Diferencias entre los videojuegos de entretenimiento y los aplicados a la educación y capacitación (Michael \& Chen, 2005).
    \item \textbf{2.3} Casos de uso de \textit{serious games} en la formación de ingeniería de software.
\end{itemize}

\noindent\textbf{3. Gamificación y Simulación}
\begin{itemize}
    \item \textbf{3.1} Concepto de gamificación en entornos universitarios de aprendizaje (Deterding et al., 2011).
    \item \textbf{3.2} Elementos y mecánicas de diseño (Misiones, puntajes y escenarios adaptativos controlados) (Sánchez, García \& Ruiz, 2021).
    \item \textbf{3.3} Efectividad de la simulación en el desarrollo de competencias y habilidades prácticas (Kapp, 2012).
    \item \textbf{3.4} Entornos virtuales inmersivos y competencias como \textit{Capture The Flag (CTF)} en ciberseguridad (Chapman et al., 2022).
\end{itemize}

\noindent\textbf{4. Tecnologías, Metodologías y Herramientas para el Desarrollo}
\begin{itemize}
    \item \textbf{4.1} Motores gráficos y lógicos de videojuegos para simuladores de entornos de red.
    \item \textbf{4.2} Diseño de interfaces de usuario (UI) y experiencia de usuario (UX) centradas en el aprendizaje.
    \item \textbf{4.3} Integración técnica de mecánicas de análisis de vulnerabilidades en código y arquitectura.
\end{itemize}